% Sections 16.1-16.7, from lectures/L16/appendix/a_can_controller.md.

\section{One node, two roles}
\label{c16:sec:roles}

A node either \textbf{transmits} its own requested frame or \textbf{receives} somebody else's, never
both for the same frame. A registered \code{role} (\code{'1'} = transmitting, \code{'0'} =
receiving), latched once per frame at the edge leaving \code{STATE_IDLE} for \code{STATE_START} and
held until the frame is over, decides which of the subblocks' ports matter. Both ways out of
\code{STATE_IDLE} require \textbf{a bus that has been idle}: with that established, \code{tx_req =
'1'} starts a transmission (\code{role = '1'}), and otherwise a dominant \code{rx_bus_s2} (somebody
else's SOF) starts a reception (\code{role = '0'}). Follow the \code{role = '1'} path in this
chapter.

\subsection*{Why ``a bus that has been idle'' does real work}
The condition earns its place twice over, for different reasons at each of the two entrances.

\paragraph{At the receive entrance}
A dominant bit by itself does not mean a frame is starting; most bits in most frames are dominant.
\code{STATE_IDLE} has to tell the dominant bit that \emph{opens} a frame from the many dominant bits
inside one, and a node in \code{STATE_IDLE} has no frame context to tell them apart with.

Two tempting rules both fail:
\begin{itemize}
  \item \textbf{A dominant level.} Fires on any dominant bit at all, so a node sitting in
    \code{STATE_IDLE} while another node transmits immediately tries to ``receive'' a frame that is
    already half over.
  \item \textbf{A recessive-to-dominant edge.} Sounds like \secref{c10:sec:frameformat}'s description
    of SOF, but bit stuffing forces a transition after every fifth identical bit, so edges in both
    directions, the recessive-to-dominant one included, occur constantly in the middle of a frame.
    The edge test fires within a few bits of the level test.
\end{itemize}

The rule that actually works is a \textbf{count of recessive bit periods}, and its threshold takes
one more step of care than the bit stuffing rule alone. Inside the stuffed region,
\secref{c10:sec:stuffing}'s rule guarantees no run longer than five identical bits. But the tail is
unstuffed and almost entirely recessive: an acknowledged frame ends with the ACK delimiter plus EOF,
eight recessive bit periods in a row, and a frame nobody acknowledges shows ten (CRC delimiter,
recessive ACK slot, ACK delimiter, EOF). A threshold of six to ten would declare the bus idle while a
frame's own tail was still passing over it, and a pending \code{tx_req} would then drive SOF out over
somebody's last EOF bits.

Count recessive bit periods in \code{STATE_IDLE} and require \textbf{eleven}, the same number real
CAN gives a node that has to judge whether the bus is idle with no frame context: the acknowledged
tail of eight plus the standard's three-bit intermission, and past even the ten-bit tail of an
unacknowledged frame. Only once the count is full does a dominant \code{rx_bus_s2} mean SOF. Two
consequences for the rest of the design:
\begin{itemize}
  \item \code{bit_timer} has to keep running in \code{STATE_IDLE}, since the count is in bit periods.
  \item Reset has to start the count \textbf{full}. A node coming out of reset has no reason to
    believe a frame is in progress, and starting it empty would make it ignore the first frame on the
    bus.
\end{itemize}

\paragraph{At the transmit entrance}
The same count answers a different question: may this node start driving the bus at all? Real CAN
says a node may only start transmitting on an idle bus, and the reason is not politeness. Arbitration
settles a dispute only between nodes that started at the \emph{same} SOF, by comparing what each sent
against what the bus reads, bit by bit, down through the identifier. A node starting in the middle of
a frame never takes part in that contest; it merely drives dominant bits over somebody else's fields.
The frame in progress is corrupted, and the damage turns up at the other end as a stuffing violation
or a CRC error, with nothing pointing back at the node that caused it. So \code{tx_req} waits on the
count too.

Together these conditions keep \code{error} observable. A node losing arbitration (\chapref{c17:ch})
falls back to \code{STATE_IDLE} while the winner is still driving the bus, and \code{STATE_START}
clears \code{error}. Under either of the two failed receive rules above it would enter
\code{STATE_START} again within a few bit periods and wipe the \code{error} it had just raised; with
the transmit entrance ungated, a caller still holding \code{tx_req} high would wipe it the very next
cycle. Gating both means the loser waits, correctly, for the next frame, and that a caller polling
\code{error} now and then still sees the error.

% ------------------------------------------------------------------------------------------
\section{The frame, field by field}
\label{c16:sec:fields}

\lecfig[0.92]{lectures/L16/appendix/images/fsm.png}{\code{can_controller}'s transmit state
machine.}{c16:fig:fsm}

Every \textbf{stuffed} field goes through \code{tx_shift_reg} as one group, sent MSB first. Every
field gets \textbf{two} states: a one-cycle \code{STATE_LOAD_*} pulsing \code{tx_shift_reg.load}, and
then the shift state waiting for \code{tx_shift_reg.done}.

\begin{booktable}{@{}p{11.6em}p{5em}r>{\raggedright\arraybackslash}p{8.6em}c@{}}
  \thead{State} & \thead{Field} & \thead{Bits} & \thead{Value sent} & \thead{Stuffed} \\
  \midrule
  \code{STATE_START} / \code{STATE_SOF} & SOF & 1 & dominant \code{'0'} & yes \\
  \code{STATE_LOAD_ARB_HI} / \code{STATE_ARB_HI} & ID high & 8 & \code{tx_id(10 downto 3)} & yes \\
  \code{STATE_LOAD_ARB_LO} / \code{STATE_ARB_LO} & ID low + RTR & 4 & \code{tx_id(2 downto 0)}, then
    \code{RTR='0'} & yes \\
  \code{STATE_LOAD_CTRL} / \code{STATE_CTRL} & Control & 6 & \code{IDE='0'}, \code{r0='0'}, then
    4-bit \code{tx_dlc} & yes \\
  \code{STATE_LOAD_DATA} / \code{STATE_DATA} & Data byte & 8 & current byte, repeated \code{DLC}
    times & yes \\
  \code{STATE_CRC_WAIT} & (settling) & 0 & let \code{crc15} settle, then latch it into \code{crc_tx}
    & no \\
  \code{STATE_LOAD_CRC_HI} / \code{STATE_CRC_HI} & CRC high & 8 & \code{crc_tx(14 downto 7)} &
    yes \\
  \code{STATE_LOAD_CRC_LO} / \code{STATE_CRC_LO} & CRC low & 7 & \code{crc_tx(6 downto 0)} & yes \\
  \code{STATE_CRC_LO_WAIT} & (settling) & 0 & let \code{crc_valid} settle & no \\
  \code{STATE_CRC_DELIM} & CRC delimiter & 1 & recessive & no \\
  \code{STATE_ACK_SLOT} & ACK slot & 1 & recessive (the receiver drives dominant) & no \\
  \code{STATE_ACK_DELIM} & ACK delimiter & 1 & recessive & no \\
  \code{STATE_EOF} & EOF & 7 & recessive & no \\
\end{booktable}

Notes on the shift register's loads:
\begin{itemize}
  \item \code{tx_shift_reg.data} is a \code{byte_t}, always 8 bits, and sends the top
    \code{bit_count} of them (\chapref{c14:ch}). A group narrower than a byte is therefore
    \textbf{left-justified and zero-padded} on load, and the padding is never sent. The table above
    names every group after its real bits; the padding is mechanical.
  \item The control byte is loaded as \code{"00" & tx_dlc & "00"} with \code{bit_count = 6}, so the
    six bits sent are the top six: \code{0}, \code{0}, \code{DLC[3:0]}.
  \item The low CRC group is loaded as \code{crc_tx(6 downto 0) & '0'} with \code{bit_count = 7}.
  \item The low ID group is loaded as \code{tx_id(2 downto 0) & '0' & "0000"} with
    \code{bit_count = 4}: the identifier's three remaining bits, and then the always dominant RTR
    that ends the arbitration field.
  \item Data selects one byte out of \code{tx_data} per \code{data_byte_idx}; loop
    \code{STATE_LOAD_DATA}/\code{STATE_DATA} once per byte, \code{DLC} times (jump straight to
    \code{STATE_CRC_WAIT} when \code{DLC = 0}).
  \item \code{tx_dlc} above 8 is outside the contract: \secref{c10:sec:frameformat}'s DLC encodes 0
    to 8, and the byte loop runs outside \code{data_byte_idx}'s range for anything larger. Real
    controllers clamp the encodings 9 to 15 to eight bytes; this design simply requires the caller not
    to send one.
\end{itemize}

The \textbf{unstuffed} closing fields (CRC delimiter, ACK slot, ACK delimiter, EOF) do not go through
\code{tx_shift_reg}; they are fixed bit counts driven directly, one bit per
\code{bit_timer.bit_done}, identically for both roles.

% ------------------------------------------------------------------------------------------
\section{What you add to the architecture}
\label{c16:sec:declare}

\file{can_controller.vhd} already holds the entity (\chapref{c11:ch}), the \code{meta_prev}
synchroniser (\chapref{c12:ch}) and all four CAN subblock instances with the signals they connect to
(\chapref{c12:ch} to \ref{c15:ch}). It holds no frame behaviour at all. Everything below is new in
this chapter.

The state type and the state machine's own signals belong in the architecture's declarative part,
above the four subblocks' signal groups:

\begin{vhdlcode}
    -- One state per frame field. Every stuffed field is a LOAD/shift pair; the
    -- trailer fields (CRC delimiter, ACK slot/delim, EOF) are driven directly.
    type state_t is (
        STATE_IDLE,
        STATE_START,
        STATE_SOF,
        STATE_LOAD_ARB_HI, STATE_ARB_HI,
        STATE_LOAD_ARB_LO, STATE_ARB_LO,
        STATE_LOAD_CTRL,   STATE_CTRL,
        STATE_LOAD_DATA,   STATE_DATA,
        STATE_CRC_WAIT,
        STATE_LOAD_CRC_HI, STATE_CRC_HI,
        STATE_LOAD_CRC_LO, STATE_CRC_LO,
        STATE_CRC_LO_WAIT,
        STATE_CRC_DELIM,
        STATE_ACK_SLOT,
        STATE_ACK_DELIM,
        STATE_EOF);

    signal state : state_t;

    -- '1' while transmitting this frame, '0' while only receiving. Latched per frame.
    signal role : std_logic;

    -- Consecutive recessive bit periods seen in STATE_IDLE, saturating at 11. Only once this is
    -- full may the node leave STATE_IDLE at all. Resets to 11, not 0.
    signal idle_bits : natural range 0 to 11;

    -- Data-byte index, and a countdown for the fixed-width trailer fields.
    signal data_byte_idx   : natural range 0 to 8;
    signal field_bit_count : natural range 0 to 7;

    -- Accumulators for the frame being received.
    signal rx_id_acc   : id_t;
    signal rx_dlc_acc  : dlc_t;
    signal rx_data_acc : data_t;

    -- This frame's own CRC, captured at STATE_CRC_WAIT before the engine is fed
    -- the CRC field's own bits and moves on.
    signal crc_tx : crc_t;

    -- The data byte selected by data_byte_idx, MSB-first (byte 0 first).
    signal tx_data_shifted : data_t;
    signal tx_current_byte : byte_t;
\end{vhdlcode}

The three \code{rx_*_acc} accumulators are not used until \chapref{c17:ch}; declaring them with the
rest keeps the group together instead of splitting it across two chapters.

Two concurrent assignments pick out the data byte currently being sent. They belong after
\code{begin}, alongside the subblock instances:

\begin{vhdlcode}
    tx_data_shifted <= std_logic_vector(shift_left(unsigned(tx_data), 8 * data_byte_idx));
    tx_current_byte <= tx_data_shifted(63 downto 56);
\end{vhdlcode}

\code{shift_left} and \code{unsigned} come from \code{numeric_std}, so add
\code{use ieee.numeric_std.all;} to the context clause if it is not already there.

Everything else in this chapter is the two processes you write: a \textbf{combinational} one driving
the subblocks per state, and a \textbf{sequential} one advancing the state machine and updating the
bookkeeping.

\paragraph{Which process drives what}
The split is worth settling before the walkthrough, since every line in it lands in one of the two
processes, and which one should be no mystery:
\begin{itemize}
  \item The \textbf{combinational} process drives every input on the subblocks (\code{bt_enable},
    \code{bt_resync}, the four \code{txsr_*} inputs, the two \code{rxsr_*} inputs in
    \chapref{c17:ch}) and the bus pair \code{tx_bus}/\code{bus_en}. Give them all a default at the
    top and override per state, in the same style as \code{bit_timer}'s own branches.
  \item \code{bt_enable} is \code{'1'} in \textbf{every} state, \code{STATE_IDLE} included, since the
    count in \code{STATE_IDLE} is measured in bit periods. The port exists so that \code{bit_timer}
    can be held, and this design never wants to.
  \item \code{txsr_shift} is \code{bt_bit_done}, but only while transmitting and only in the
    \textbf{shift states}, never in a \code{STATE_LOAD_*} state. That is what keeps
    \chapref{c14:ch}'s promise that ``load and shift are mutually exclusive'': a load state issues a
    load and nothing else.
  \item The \textbf{sequential} process owns \code{state} and the bookkeeping (\code{role},
    \code{idle_bits}, \code{data_byte_idx}, \code{field_bit_count}, \code{crc_tx}, the
    \code{rx_*_acc} accumulators) and every register-side output port: \code{tx_done},
    \code{rx_valid}, \code{error} and the \code{rx_*} registers. The bus pair stays combinational,
    driven per state as above. \code{tx_done} and \code{rx_valid} are one-cycle pulses, so default
    them to \code{'0'} at every edge and raise them where they belong.
  \item \code{error} is cleared by \code{reset_s2_n} as well as by \code{STATE_START}. It is a level a
    caller polls, so it has to read \code{'0'} out of reset rather than whatever a flip-flop happened
    to power up as.
  \item \code{tx_id}, \code{tx_dlc} and \code{tx_data} are \textbf{not} latched at
    \code{STATE_START}; the load states read them live, each at the moment its group is loaded. That
    is a deliberate simplification and it puts the obligation on the caller: hold all three stable
    from \code{tx_req} to \code{tx_done}. A register block in the same clock domain does that for
    free, and that is precisely the caller \secref{c11:sec:interface} assumes.
\end{itemize}

% ------------------------------------------------------------------------------------------
\section{The combinational process}
\label{c16:sec:comb}

It drives the subblocks' inputs and the bus pair from \code{state}, with the defaults first and the
per-state overrides after, and it is compact enough to show whole. Three things to see in it before
the code. Every \code{STATE_LOAD_*} arm loads its group and shares an arm with the paired shift
state, since the pair also shares a \code{bit_count}. \code{txsr_shift} is \code{bt_bit_done} only in
the shift states, never in a \code{STATE_LOAD_*} state: \code{load} already presents the new group's
first bit during the load cycle (\chapref{c14:ch}), so a dedicated one-cycle load state is what keeps
that promise; getting this the wrong way round is the frame-level version of the one-bit timing error
\chapref{c14:ch} warned about. And the bus is open-drain: a node drives \emph{actively} only when the
bit is dominant, and otherwise releases, so that another node can pull the same wire dominant during
the same period.

\begin{vhdlcode}
    COMBINATIONAL_PROCESS: process(state, role, bt_bit_done, tx_id, tx_dlc,
                                   tx_current_byte, crc_tx, txsr_tx_bit) is
    begin
        -- Defaults: every sub-block input, every cycle, then per-state overrides.
        bt_enable      <= '1';            -- Never held; the idle count needs it.
        bt_resync      <= '0';
        txsr_load      <= '0';
        txsr_data      <= (others => '0');
        txsr_bit_count <= "0000";
        txsr_shift     <= '0';
        rxsr_enable    <= '0';            -- The receive side arrives in L17.
        rxsr_bit_count <= "0000";
        tx_bus         <= '1';            -- Release the bus unless driving dominant.
        bus_en         <= '0';

        -- Each STATE_LOAD_* arm loads its chunk. The paired shifting state keeps
        -- the same bit_count applied, so grouping the pair in one arm is natural.
        case state is
            when STATE_START =>
                bt_resync <= '1';         -- Adopt SOF as this node's bit boundary.
                if (role = '1') then
                    txsr_load      <= '1';
                    txsr_data      <= "00000000";               -- SOF: one dominant bit.
                    txsr_bit_count <= "0001";
                end if;
            when STATE_LOAD_ARB_HI | STATE_ARB_HI =>
                if ((role = '1') and (state = STATE_LOAD_ARB_HI)) then
                    txsr_load <= '1';
                    txsr_data <= tx_id(10 downto 3);
                end if;
                txsr_bit_count <= "1000";
            when STATE_LOAD_ARB_LO | STATE_ARB_LO =>
                if ((role = '1') and (state = STATE_LOAD_ARB_LO)) then
                    txsr_load <= '1';
                    txsr_data <= tx_id(2 downto 0) & '0' & "0000";  -- ID low, then RTR.
                end if;
                txsr_bit_count <= "0100";
            when STATE_LOAD_CTRL | STATE_CTRL =>
                if ((role = '1') and (state = STATE_LOAD_CTRL)) then
                    txsr_load <= '1';
                    txsr_data <= "00" & tx_dlc & "00";              -- IDE, r0, DLC.
                end if;
                txsr_bit_count <= "0110";
            when STATE_LOAD_DATA | STATE_DATA =>
                if ((role = '1') and (state = STATE_LOAD_DATA)) then
                    txsr_load <= '1';
                    txsr_data <= tx_current_byte;
                end if;
                txsr_bit_count <= "1000";
            when STATE_LOAD_CRC_HI | STATE_CRC_HI =>
                if ((role = '1') and (state = STATE_LOAD_CRC_HI)) then
                    txsr_load <= '1';
                    txsr_data <= crc_tx(14 downto 7);
                end if;
                txsr_bit_count <= "1000";
            when STATE_LOAD_CRC_LO | STATE_CRC_LO =>
                if ((role = '1') and (state = STATE_LOAD_CRC_LO)) then
                    txsr_load <= '1';
                    txsr_data <= crc_tx(6 downto 0) & '0';
                end if;
                txsr_bit_count <= "0111";
            when others =>
                null;
        end case;

        -- Shifting states advance the register once per bit period; a LOAD
        -- state never does, which keeps load and shift mutually exclusive.
        if (role = '1') then
            case state is
                when STATE_SOF | STATE_ARB_HI | STATE_ARB_LO | STATE_CTRL
                   | STATE_DATA | STATE_CRC_HI | STATE_CRC_LO =>
                    txsr_shift <= bt_bit_done;
                when others =>
                    null;
            end case;
        end if;

        -- Driving the bus, open-drain: only the owning states, only dominant
        -- bits. A receiver's one dominant bit is the ACK slot (L17).
        case state is
            when STATE_SOF
               | STATE_LOAD_ARB_HI | STATE_ARB_HI
               | STATE_LOAD_ARB_LO | STATE_ARB_LO
               | STATE_LOAD_CTRL   | STATE_CTRL
               | STATE_LOAD_DATA   | STATE_DATA
               | STATE_LOAD_CRC_HI | STATE_CRC_HI
               | STATE_LOAD_CRC_LO | STATE_CRC_LO =>
                if ((role = '1') and (txsr_tx_bit = '0')) then
                    bus_en <= '1';
                    tx_bus <= '0';
                end if;
            when others =>
                null;
        end case;
    end process;
\end{vhdlcode}

\begin{warning}
The bounding of the bus driving in that last \code{case} statement deserves a sentence of its own:
\code{txsr_tx_bit} is a level and holds its most recent value forever, so an unbounded version goes
on driving after the group that produced it. A frame whose CRC ends on a dominant bit would still
pull the bus down through the whole tail, and a node sitting in \code{STATE_IDLE} with \code{role}
still \code{'1'} from its last frame would hold the entire bus dominant forever. The tail's fields
are recessive for both roles and drive nothing, with the single exception of a receiver's ACK slot
(\chapref{c17:ch}).
\end{warning}

% ------------------------------------------------------------------------------------------
\section{The sequential process}
\label{c16:sec:seq}

It owns \code{state}, the bookkeeping and the output ports. Skeleton, reset and defaults first:

\begin{vhdlcode}
    SEQUENTIAL_PROCESS: process(clock, reset_s2_n) is
    begin
        if (reset_s2_n = '0') then
            state           <= STATE_IDLE;
            role            <= '0';
            idle_bits       <= 11;        -- Full: an idle bus is assumed at power-up.
            data_byte_idx   <= 0;
            field_bit_count <= 0;
            rx_id_acc       <= (others => '0');
            rx_dlc_acc      <= (others => '0');
            rx_data_acc     <= (others => '0');
            crc_tx          <= (others => '0');
            tx_done         <= '0';
            rx_valid        <= '0';
            error           <= '0';
            rx_id           <= (others => '0');
            rx_dlc          <= (others => '0');
            rx_data         <= (others => '0');
        elsif (rising_edge(clock)) then
            tx_done  <= '0';              -- The two pulse outputs default low;
            rx_valid <= '0';              -- error is a level and holds.

            case state is
\end{vhdlcode}

Everything below is arms of this single \code{case} statement. First the two states every frame
passes through once:

\begin{vhdlcode}
                when STATE_IDLE =>
                    if (bt_sample = '1') then
                        if (rx_bus_s2 = '1') then
                            if (idle_bits /= 11) then
                                idle_bits <= idle_bits + 1;
                            end if;
                        else
                            idle_bits <= 0;
                        end if;
                    end if;
                    if (idle_bits = 11) then
                        if (tx_req = '1') then
                            role      <= '1';
                            idle_bits <= 0;
                            state     <= STATE_START;
                        elsif (rx_bus_s2 = '0') then
                            role      <= '0';    -- Someone else's SOF: receive (L17).
                            idle_bits <= 0;
                            state     <= STATE_START;
                        end if;
                    end if;
                when STATE_START =>
                    error         <= '0';
                    data_byte_idx <= 0;
                    rx_data_acc   <= (others => '0');
                    state         <= STATE_SOF;
\end{vhdlcode}

\code{STATE_IDLE} counts recessive bit periods at the sample point and saturates at 11; a dominant
sample clears the count. Nothing leaves the state until the count is \textbf{full}, whatever
\code{tx_req} and \code{rx_bus_s2} do, so a \code{tx_req} pulse arriving while the bus is busy is
\textbf{discarded}, not queued; the caller asks again. \code{STATE_START} needs no code for the SOF
load or the resync, since the combinational process issues both from the state itself; setup that has
to be \emph{remembered}, clearing \code{error}, the byte index and the received-data accumulator, is
what lands here. \code{rx_data_acc} needs the clear because a received frame writes only the bytes its
DLC covers (\chapref{c17:ch}): a five-byte frame arriving after an eight-byte frame has to report
zeros in bytes 5 to 7, not the previous frame's residue. \code{rx_id_acc} and \code{rx_dlc_acc} are
always written in full, so they need none. (\code{crc15} is cleared by this state too; see
\secref{c16:sec:crcclear}.)

Then the LOAD/shift pairs. This is the shape every stuffed field follows, shown once for the
identifier's high group; the transmit half of a LOAD state moves on after its single cycle, and the
shift state waits for \code{tx_shift_reg.done}.

\begin{aside}
A timing note before the waveforms surprise you: \chapref{c14:ch} described \code{load} as issued at
the very moment of the previous group's last \code{bit_done}, and here it is issued two clock edges
later, one for the state machine to see \code{txsr_done} and one for the LOAD state itself. Every
group's first bit therefore reaches the wire about 40~ns into its 1000~ns bit period; the sample
point at 70~\% absorbs that with room to spare, and the offset does not accumulate, since every group
restarts from the same \code{bit_done}.
\end{aside}

\begin{vhdlcode}
                when STATE_SOF =>
                    if (role = '1') then
                        if (txsr_done = '1') then
                            state <= STATE_LOAD_ARB_HI;
                        end if;
                    end if;                       -- role = '0': L17.
                when STATE_LOAD_ARB_HI =>
                    if (role = '1') then
                        state <= STATE_ARB_HI;    -- One cycle; the load is issued.
                    end if;                       -- role = '0': L17.
                when STATE_ARB_HI =>
                    if (role = '1') then
                        if (txsr_done = '1') then -- Arbitration watching joins here in L17.
                            state <= STATE_LOAD_ARB_LO;
                        end if;
                    end if;                       -- role = '0': L17.
                when STATE_LOAD_ARB_LO =>
                    if (role = '1') then
                        state <= STATE_ARB_LO;
                    end if;                       -- role = '0': L17.
                when STATE_ARB_LO =>
                    if (role = '1') then
                        if (txsr_done = '1') then -- Also watched in L17; RTR is dominant,
                            state <= STATE_LOAD_CTRL;  -- so it can never be the bit that loses.
                        end if;
                    end if;                       -- role = '0': L17.
\end{vhdlcode}

The control and data states have the same shape with two decisions attached: the control field's
\code{done} consults the DLC, and the data pair loops once per byte:

\begin{vhdlcode}
                when STATE_LOAD_CTRL =>
                    if (role = '1') then
                        state <= STATE_CTRL;
                    end if;                       -- role = '0': L17.
                when STATE_CTRL =>
                    if (role = '1') then
                        if (txsr_done = '1') then
                            if (unsigned(tx_dlc) = 0) then
                                state <= STATE_CRC_WAIT;    -- No data field at all.
                            else
                                state <= STATE_LOAD_DATA;
                            end if;
                        end if;
                    end if;                       -- role = '0': L17.
                when STATE_LOAD_DATA =>
                    if (role = '1') then
                        state <= STATE_DATA;
                    end if;                       -- role = '0': L17.
                when STATE_DATA =>
                    if (role = '1') then
                        if (txsr_done = '1') then
                            if (data_byte_idx + 1 = to_integer(unsigned(tx_dlc))) then
                                state <= STATE_CRC_WAIT;
                            else
                                data_byte_idx <= data_byte_idx + 1;
                                state         <= STATE_LOAD_DATA;
                            end if;
                        end if;
                    end if;                       -- role = '0': L17.
\end{vhdlcode}

Then the CRC. \code{STATE_CRC_WAIT} is an empty cycle so that the last data bit has time to finish
being integrated before \code{crc_value} is read, and its real job is the latching: both CRC groups
are loaded from \code{crc_tx}, never from the live engine (\secref{c16:sec:crccapture} explains what
goes wrong otherwise). The two CRC pairs have the arbitration pairs' shape again:

\begin{vhdlcode}
                when STATE_CRC_WAIT =>
                    crc_tx <= crc_value;          -- Capture before the engine moves on.
                    state  <= STATE_LOAD_CRC_HI;
                when STATE_LOAD_CRC_HI =>
                    if (role = '1') then
                        state <= STATE_CRC_HI;
                    end if;                       -- role = '0': L17.
                when STATE_CRC_HI =>
                    if (role = '1') then
                        if (txsr_done = '1') then
                            state <= STATE_LOAD_CRC_LO;
                        end if;
                    end if;                       -- role = '0': L17.
                when STATE_LOAD_CRC_LO =>
                    if (role = '1') then
                        state <= STATE_CRC_LO;
                    end if;                       -- role = '0': L17.
                when STATE_CRC_LO =>
                    if (role = '1') then
                        if (txsr_done = '1') then
                            state <= STATE_CRC_LO_WAIT;
                        end if;
                    end if;                       -- role = '0': L17.
\end{vhdlcode}

Finally the tail, where the shift register bows out and the state machine counts \code{bt_bit_done}
pulses itself. \code{STATE_CRC_LO_WAIT} is one clock cycle for a transmitter, \textbf{not} a bit
period: \code{tx_shift_reg.done} pulses at the closing shift, which is itself a \code{bit_done}, so
the transmitter is already on a bit boundary and the tail can begin immediately; waiting a whole
period would push an extra recessive bit out onto the wire. (A receiver, \emph{by contrast}, waits a
whole \code{bit_done} here, and checks \code{crc_valid} while it waits; \chapref{c17:ch}.)

\begin{vhdlcode}
                when STATE_CRC_LO_WAIT =>
                    if (role = '1') then
                        state <= STATE_CRC_DELIM; -- One clock cycle, not a bit period.
                    end if;                       -- role = '0': L17 checks crc_valid here.
                when STATE_CRC_DELIM =>
                    if (bt_bit_done = '1') then
                        state <= STATE_ACK_SLOT;
                    end if;
                when STATE_ACK_SLOT =>
                    if (bt_bit_done = '1') then   -- Transmitter releases; a receiver
                        state <= STATE_ACK_DELIM; -- acknowledges here (L17).
                    end if;
                when STATE_ACK_DELIM =>
                    if (bt_bit_done = '1') then
                        field_bit_count <= 6;     -- Arm the 7-bit EOF countdown.
                        state           <= STATE_EOF;
                    end if;
                when STATE_EOF =>
                    if (bt_bit_done = '1') then
                        if (field_bit_count = 0) then
                            if (role = '1') then
                                tx_done <= '1';
                            end if;                -- role = '0': L17 latches rx_* here.
                            state <= STATE_IDLE;
                        else
                            field_bit_count <= field_bit_count - 1;
                        end if;
                    end if;
            end case;
        end if;
    end process;
\end{vhdlcode}

The tail's states need no \code{role} test to move on, since both roles count the same whole bit
periods through the same fields; only what happens \emph{inside} them differs, and those differences
are all \chapref{c17:ch}'s.

% ------------------------------------------------------------------------------------------
\section{Capturing the CRC}
\label{c16:sec:crccapture}

\code{STATE_CRC_WAIT} latches \code{crc_value} into \code{crc_tx}, and both CRC groups are loaded
from \code{crc_tx} rather than from the engine. That detour looks redundant and is not, for a reason
worth following closely, since it is the one place in this design where two individually correct
rules together make a bug.

Rule one: the engine keeps running through the CRC field, since that is what returns it to zero
(\secref{c16:sec:crcfeed}). Rule two: every group is loaded when its \code{STATE_LOAD_*} state runs.

Now put them together. \code{STATE_LOAD_CRC_HI} runs first and reads the right value. Then
\code{STATE_CRC_HI} shifts out eight bits, and every one of them goes back into the engine, which
dutifully advances the register. By the time \code{STATE_LOAD_CRC_LO} runs, \code{crc_value} is
\textbf{no longer the frame's CRC}: it is the register eight bits further on, which for a correct
frame means the remaining seven CRC bits shifted up with zeros shifted in beneath them. Load
\code{crc_value(6 downto 0)} there, and the frame's last seven CRC bits go out as seven zeros, in
every frame, always.

The failure is silent in the worst possible way. The transmitter is perfectly content; nothing it can
see is wrong. The receiver's engine does not reach zero, so it reports a CRC error, and the obvious
place to look is the receive path, which is entirely reasonable. Latching the value once, at the
moment it is finished and before anything else touches the engine, removes the interaction entirely.

% ------------------------------------------------------------------------------------------
\section{Feeding \code{crc15}}
\label{c16:sec:crcfeed}

\code{crc15} has to see every real bit from SOF through the end of the \textbf{CRC field}, and no
others. Gate its enable on \code{tx_shift_reg.bit_valid} (\code{'0'} at every group's closing shift,
so that the unchanged bit is not fed in twice), and shut out the stuff bits and the receive path:

\begin{vhdlcode}
crc_enable  <= txsr_bit_valid and not txsr_stuff and role;
crc_data_in <= txsr_tx_bit;
\end{vhdlcode}

Three things follow from this being a simple concurrent assignment with no state term in it.

\textbf{The CRC field is fed back into the engine, and that is deliberate.} The value being sent sits
safely in \code{crc_tx}, so nothing has to read the engine again once \code{STATE_CRC_WAIT} has
captured it, and the expression has no state term that would stop the feed at the data field. On the
transmit side the feedback is merely harmless: \chapref{c13:ch}'s generate-then-check property means
that a message followed by its own CRC returns the register to zero, and \code{STATE_START}'s clear
(\secref{c16:sec:crcclear}) hands the next frame a fresh engine anyway. Where the feedback earns its
place is on the \textbf{receive side}: \chapref{c17:ch}'s CRC check \emph{is} the observation that the
register has returned to zero once the received CRC field has gone in, and that observation exists
only because the CRC field is fed back. One rule for both roles, every real bit from SOF through the
end of the CRC field, is what lets a single engine generate on one node while it checks on another.

\textbf{The tail excludes itself.} CRC delimiter, ACK slot, ACK delimiter and EOF never go through
\code{tx_shift_reg} at all, so \code{txsr_bit_valid} is already low across all of them. No state term
is needed to keep them out of the checksum.

\textbf{\code{role} is the placeholder for the receive side.} With \code{role = '0'} this expression
holds \code{crc15} disabled, which is correct for this chapter (transmission only) but not for the
finished controller. \Chapref{c17:ch} replaces the whole assignment with one feeding
\code{rx_shift_reg.real_bit} on \code{rx_shift_reg.real_bit_valid} during reception; the transmit half
above is unchanged by that.

% ------------------------------------------------------------------------------------------
\section{Clearing \code{crc15}}
\label{c16:sec:crcclear}

One more concurrent assignment, and it is the one that makes the engine safe to reuse frame after
frame:

\begin{vhdlcode}
crc_clear <= '1' when state = STATE_START else '0';
\end{vhdlcode}

On the happy path it does nothing at all, since a frame that runs all the way has already returned
the register to zero. It is there for the frames that do \emph{not} finish. \Chapref{c17:ch} gives
this node three ways to abandon a frame halfway (it loses arbitration, it sees a stuffing violation,
its CRC check fails), and each leaves the engine holding a half-finished value with no prospect of
ever reaching zero by itself. Every frame the node sent after that would carry a CRC computed from
that residue, and every frame it received would fail the check, until somebody pressed reset.

Two things are worth noticing. \code{STATE_START} is the right home for it since it is already the
``set this frame up'' state, alongside clearing \code{error} and the byte index; the clear belongs
there rather than among the CRC states. And \code{can_controller_tb} cannot catch a missing clear,
since neither of its two scenarios lets a node transmit again after an abort. That makes this a rule
to get right by reasoning rather than by waiting for a red testbench, the same lesson \code{crc_tx}
gave in \secref{c16:sec:crccapture}: the silent failures are the ones to reason out in advance.
